% Options for packages loaded elsewhere
% Options for packages loaded elsewhere
\PassOptionsToPackage{unicode}{hyperref}
\PassOptionsToPackage{hyphens}{url}
%
\documentclass[
  british,
  a4paper,
]{scrbook}
\usepackage{xcolor}
\usepackage[margin=1in]{geometry}
\usepackage{amsmath,amssymb}
\setcounter{secnumdepth}{5}
\usepackage{iftex}
\ifPDFTeX
  \usepackage[T1]{fontenc}
  \usepackage[utf8]{inputenc}
  \usepackage{textcomp} % provide euro and other symbols
\else % if luatex or xetex
  \usepackage{unicode-math} % this also loads fontspec
  \defaultfontfeatures{Scale=MatchLowercase}
  \defaultfontfeatures[\rmfamily]{Ligatures=TeX,Scale=1}
\fi
\usepackage{lmodern}
\ifPDFTeX\else
  % xetex/luatex font selection
\fi
% Use upquote if available, for straight quotes in verbatim environments
\IfFileExists{upquote.sty}{\usepackage{upquote}}{}
\IfFileExists{microtype.sty}{% use microtype if available
  \usepackage[]{microtype}
  \UseMicrotypeSet[protrusion]{basicmath} % disable protrusion for tt fonts
}{}
\makeatletter
\@ifundefined{KOMAClassName}{% if non-KOMA class
  \IfFileExists{parskip.sty}{%
    \usepackage{parskip}
  }{% else
    \setlength{\parindent}{0pt}
    \setlength{\parskip}{6pt plus 2pt minus 1pt}}
}{% if KOMA class
  \KOMAoptions{parskip=half}}
\makeatother
% Make \paragraph and \subparagraph free-standing
\makeatletter
\ifx\paragraph\undefined\else
  \let\oldparagraph\paragraph
  \renewcommand{\paragraph}{
    \@ifstar
      \xxxParagraphStar
      \xxxParagraphNoStar
  }
  \newcommand{\xxxParagraphStar}[1]{\oldparagraph*{#1}\mbox{}}
  \newcommand{\xxxParagraphNoStar}[1]{\oldparagraph{#1}\mbox{}}
\fi
\ifx\subparagraph\undefined\else
  \let\oldsubparagraph\subparagraph
  \renewcommand{\subparagraph}{
    \@ifstar
      \xxxSubParagraphStar
      \xxxSubParagraphNoStar
  }
  \newcommand{\xxxSubParagraphStar}[1]{\oldsubparagraph*{#1}\mbox{}}
  \newcommand{\xxxSubParagraphNoStar}[1]{\oldsubparagraph{#1}\mbox{}}
\fi
\makeatother

\usepackage{color}
\usepackage{fancyvrb}
\newcommand{\VerbBar}{|}
\newcommand{\VERB}{\Verb[commandchars=\\\{\}]}
\DefineVerbatimEnvironment{Highlighting}{Verbatim}{commandchars=\\\{\}}
% Add ',fontsize=\small' for more characters per line
\usepackage{framed}
\definecolor{shadecolor}{RGB}{241,243,245}
\newenvironment{Shaded}{\begin{snugshade}}{\end{snugshade}}
\newcommand{\AlertTok}[1]{\textcolor[rgb]{0.68,0.00,0.00}{#1}}
\newcommand{\AnnotationTok}[1]{\textcolor[rgb]{0.37,0.37,0.37}{#1}}
\newcommand{\AttributeTok}[1]{\textcolor[rgb]{0.40,0.45,0.13}{#1}}
\newcommand{\BaseNTok}[1]{\textcolor[rgb]{0.68,0.00,0.00}{#1}}
\newcommand{\BuiltInTok}[1]{\textcolor[rgb]{0.00,0.23,0.31}{#1}}
\newcommand{\CharTok}[1]{\textcolor[rgb]{0.13,0.47,0.30}{#1}}
\newcommand{\CommentTok}[1]{\textcolor[rgb]{0.37,0.37,0.37}{#1}}
\newcommand{\CommentVarTok}[1]{\textcolor[rgb]{0.37,0.37,0.37}{\textit{#1}}}
\newcommand{\ConstantTok}[1]{\textcolor[rgb]{0.56,0.35,0.01}{#1}}
\newcommand{\ControlFlowTok}[1]{\textcolor[rgb]{0.00,0.23,0.31}{\textbf{#1}}}
\newcommand{\DataTypeTok}[1]{\textcolor[rgb]{0.68,0.00,0.00}{#1}}
\newcommand{\DecValTok}[1]{\textcolor[rgb]{0.68,0.00,0.00}{#1}}
\newcommand{\DocumentationTok}[1]{\textcolor[rgb]{0.37,0.37,0.37}{\textit{#1}}}
\newcommand{\ErrorTok}[1]{\textcolor[rgb]{0.68,0.00,0.00}{#1}}
\newcommand{\ExtensionTok}[1]{\textcolor[rgb]{0.00,0.23,0.31}{#1}}
\newcommand{\FloatTok}[1]{\textcolor[rgb]{0.68,0.00,0.00}{#1}}
\newcommand{\FunctionTok}[1]{\textcolor[rgb]{0.28,0.35,0.67}{#1}}
\newcommand{\ImportTok}[1]{\textcolor[rgb]{0.00,0.46,0.62}{#1}}
\newcommand{\InformationTok}[1]{\textcolor[rgb]{0.37,0.37,0.37}{#1}}
\newcommand{\KeywordTok}[1]{\textcolor[rgb]{0.00,0.23,0.31}{\textbf{#1}}}
\newcommand{\NormalTok}[1]{\textcolor[rgb]{0.00,0.23,0.31}{#1}}
\newcommand{\OperatorTok}[1]{\textcolor[rgb]{0.37,0.37,0.37}{#1}}
\newcommand{\OtherTok}[1]{\textcolor[rgb]{0.00,0.23,0.31}{#1}}
\newcommand{\PreprocessorTok}[1]{\textcolor[rgb]{0.68,0.00,0.00}{#1}}
\newcommand{\RegionMarkerTok}[1]{\textcolor[rgb]{0.00,0.23,0.31}{#1}}
\newcommand{\SpecialCharTok}[1]{\textcolor[rgb]{0.37,0.37,0.37}{#1}}
\newcommand{\SpecialStringTok}[1]{\textcolor[rgb]{0.13,0.47,0.30}{#1}}
\newcommand{\StringTok}[1]{\textcolor[rgb]{0.13,0.47,0.30}{#1}}
\newcommand{\VariableTok}[1]{\textcolor[rgb]{0.07,0.07,0.07}{#1}}
\newcommand{\VerbatimStringTok}[1]{\textcolor[rgb]{0.13,0.47,0.30}{#1}}
\newcommand{\WarningTok}[1]{\textcolor[rgb]{0.37,0.37,0.37}{\textit{#1}}}

\usepackage{longtable,booktabs,array}
\newcounter{none} % for unnumbered tables
\usepackage{calc} % for calculating minipage widths
% Correct order of tables after \paragraph or \subparagraph
\usepackage{etoolbox}
\makeatletter
\patchcmd\longtable{\par}{\if@noskipsec\mbox{}\fi\par}{}{}
\makeatother
% Allow footnotes in longtable head/foot
\IfFileExists{footnotehyper.sty}{\usepackage{footnotehyper}}{\usepackage{footnote}}
\makesavenoteenv{longtable}
\usepackage{graphicx}
\makeatletter
\newsavebox\pandoc@box
\newcommand*\pandocbounded[1]{% scales image to fit in text height/width
  \sbox\pandoc@box{#1}%
  \Gscale@div\@tempa{\textheight}{\dimexpr\ht\pandoc@box+\dp\pandoc@box\relax}%
  \Gscale@div\@tempb{\linewidth}{\wd\pandoc@box}%
  \ifdim\@tempb\p@<\@tempa\p@\let\@tempa\@tempb\fi% select the smaller of both
  \ifdim\@tempa\p@<\p@\scalebox{\@tempa}{\usebox\pandoc@box}%
  \else\usebox{\pandoc@box}%
  \fi%
}
% Set default figure placement to htbp
\def\fps@figure{htbp}
\makeatother



\ifLuaTeX
\usepackage[bidi=basic,shorthands=off]{babel}
\else
\usepackage[bidi=default,shorthands=off]{babel}
\fi
\ifLuaTeX
  \usepackage{selnolig} % disable illegal ligatures
\fi


\setlength{\emergencystretch}{3em} % prevent overfull lines

\providecommand{\tightlist}{%
  \setlength{\itemsep}{0pt}\setlength{\parskip}{0pt}}



 


\makeatletter
\@ifpackageloaded{tcolorbox}{}{\usepackage[skins,breakable]{tcolorbox}}
\@ifpackageloaded{fontawesome5}{}{\usepackage{fontawesome5}}
\definecolor{quarto-callout-color}{HTML}{909090}
\definecolor{quarto-callout-note-color}{HTML}{0758E5}
\definecolor{quarto-callout-important-color}{HTML}{CC1914}
\definecolor{quarto-callout-warning-color}{HTML}{EB9113}
\definecolor{quarto-callout-tip-color}{HTML}{00A047}
\definecolor{quarto-callout-caution-color}{HTML}{FC5300}
\definecolor{quarto-callout-color-frame}{HTML}{acacac}
\definecolor{quarto-callout-note-color-frame}{HTML}{4582ec}
\definecolor{quarto-callout-important-color-frame}{HTML}{d9534f}
\definecolor{quarto-callout-warning-color-frame}{HTML}{f0ad4e}
\definecolor{quarto-callout-tip-color-frame}{HTML}{02b875}
\definecolor{quarto-callout-caution-color-frame}{HTML}{fd7e14}
\makeatother
\makeatletter
\@ifpackageloaded{bookmark}{}{\usepackage{bookmark}}
\makeatother
\makeatletter
\@ifpackageloaded{caption}{}{\usepackage{caption}}
\AtBeginDocument{%
\ifdefined\contentsname
  \renewcommand*\contentsname{Table of contents}
\else
  \newcommand\contentsname{Table of contents}
\fi
\ifdefined\listfigurename
  \renewcommand*\listfigurename{List of Figures}
\else
  \newcommand\listfigurename{List of Figures}
\fi
\ifdefined\listtablename
  \renewcommand*\listtablename{List of Tables}
\else
  \newcommand\listtablename{List of Tables}
\fi
\ifdefined\figurename
  \renewcommand*\figurename{Figure}
\else
  \newcommand\figurename{Figure}
\fi
\ifdefined\tablename
  \renewcommand*\tablename{Table}
\else
  \newcommand\tablename{Table}
\fi
}
\@ifpackageloaded{float}{}{\usepackage{float}}
\floatstyle{ruled}
\@ifundefined{c@chapter}{\newfloat{codelisting}{h}{lop}}{\newfloat{codelisting}{h}{lop}[chapter]}
\floatname{codelisting}{Listing}
\newcommand*\listoflistings{\listof{codelisting}{List of Listings}}
\makeatother
\makeatletter
\makeatother
\makeatletter
\@ifpackageloaded{caption}{}{\usepackage{caption}}
\@ifpackageloaded{subcaption}{}{\usepackage{subcaption}}
\makeatother
\usepackage{bookmark}
\IfFileExists{xurl.sty}{\usepackage{xurl}}{} % add URL line breaks if available
\urlstyle{same}
\hypersetup{
  pdftitle={OtsuFire: an R tutorial},
  pdfauthor={Natalia Quintero, Olga Viedma, Jose Manuel Moreno},
  pdflang={en-GB},
  hidelinks,
  pdfcreator={LaTeX via pandoc}}


\title{OtsuFire: an R tutorial}
\usepackage{etoolbox}
\makeatletter
\providecommand{\subtitle}[1]{% add subtitle to \maketitle
  \apptocmd{\@title}{\par {\large #1 \par}}{}{}
}
\makeatother
\subtitle{Annual burned-area mapping with an interpretable two-stage
framework}
\author{Natalia Quintero, Olga Viedma, Jose Manuel Moreno}
\date{2026-06-01}
\begin{document}
\frontmatter
\maketitle

\renewcommand*\contentsname{Table of contents}
{
\setcounter{tocdepth}{2}
\tableofcontents
}

\mainmatter
\bookmarksetup{startatroot}

\chapter*{Welcome}\label{welcome}
\addcontentsline{toc}{chapter}{Welcome}

\markboth{Welcome}{Welcome}

This is a tutorial for the \textbf{OtsuFire} R package: an open-source
framework for annual burned-area mapping from medium-resolution optical
imagery, built around a hybrid, intrinsically interpretable design.

The tutorial is organised as a progressive walkthrough of the framework.
The first chapter describes OtsuFire conceptually: what it is, what gap
it fills, how the workflow is organised in two stages, and what input
information it relies on. The following chapters then open up the
accompanying R package, illustrating each stage of the framework with
worked examples that the reader can run on their own data.

The book is written as a living document. Chapters are added and revised
as the tutorial grows, so its structure reflects the current state of
the OtsuFire framework rather than a fixed reference.

\begin{tcolorbox}[enhanced jigsaw, arc=.35mm, bottomrule=.15mm, bottomtitle=1mm, breakable, colback=white, colbacktitle=quarto-callout-note-color!10!white, colframe=quarto-callout-note-color-frame, coltitle=black, left=2mm, leftrule=.75mm, opacityback=0, opacitybacktitle=0.6, rightrule=.15mm, title=\textcolor{quarto-callout-note-color}{\faInfo}\hspace{0.5em}{Note}, titlerule=0mm, toprule=.15mm, toptitle=1mm]

The package is distributed through CRAN as \texttt{OtsuFire}.
Installation instructions and the first worked example are provided in
the \emph{How to use the OtsuFire package} chapter.

\end{tcolorbox}

\bookmarksetup{startatroot}

\chapter{OtsuFire}\label{otsufire}

\section{Introduction}\label{introduction}

This tutorial introduces the OtsuFire framework through the accompanying
R package of the same name. The package provides a reference
implementation of the framework, and the tutorial uses it as a practical
entry point: each conceptual step of the workflow is explained alongside
the package function that realises it, so that users learn the method
and the tool together.

OtsuFire targets annual burned-area mapping from medium-resolution
optical imagery in settings where externally compiled training labels
are sparse, inconsistent or unavailable. Its design priority is
interpretability. Every candidate burned patch produced by the workflow
carries the chain of decisions that led to its classification, which
allows users to trace, audit and reproduce each mapping outcome. The
framework is therefore positioned within the family of intrinsically
interpretable GeoAI approaches, in which interpretability is built into
the mapping architecture rather than added afterwards through post-hoc
explanations.

The following sections describe the framework conceptually and summarise
the input information it relies on. The remainder of the tutorial then
opens up the package step by step, illustrating how each component of
the framework is configured and executed in practice.

\section{OtsuFire}\label{otsufire-1}

OtsuFire is a hybrid, two-stage framework for annual burned-area
mapping. It is built on the premise that burned-area cartography
benefits from combining an unsupervised stage with a transparent
rule-based core with a supervised stage layer, rather than from either
approach in isolation. The rule-based core delineates a candidate-patch
universe from spectral-change information using Otsu thresholding and
explicit spatial rules. The supervised stage, when applied, refines this
candidate universe by assigning a burned probability to each patch using
a supervised classifier trained on labels derived from the rule-based
core itself.

The two stages are hierarchical rather than interchangeable. The
unsupervised phase produces the candidate universe and the auditable
evidence supporting each polygon-level decision. The supervised phase
operates on that same universe, re-ranking patches according to their
probability of being burned, but never redefining the set of candidates.
As a result, the supervised layer remains dependent on, and constrained
by, the spatial decisions of the unsupervised stage. This separation
allows the framework to combine the traceability of explicit rules with
the discriminative power of supervised learning, without inheriting the
full dependence of conventional supervised pipelines on externally
compiled training data.

The framework was developed and calibrated for the Iberian Peninsula
using the Relativised Burn Ratio (RBR) as the primary spectral-change
index, but its architecture is index-agnostic and transferable.
Application to other regions, sensors or spectral-change metrics
requires recalibration of vegetation stratum-specific thresholds, growth
parameters and lower bounds, but does not require changes to the
structure of the workflow.

\subsection{Unsupervised stage}\label{unsupervised-stage}

The unsupervised stage transforms an annual spectral-change composite
into an auditable layer of candidate burned patches, each assigned to
one of three classes that summarise the strength of the evidence
supporting it.

The stage begins with stratified Otsu thresholding. The RBR distribution
is partitioned into strata defined by the intersection of vegetation
classes and ecoregions, so that thresholds reflect the local background
reflectance, land cover and fire effects rather than a global histogram.
For each stratum, Otsu's method selects the threshold that maximises
between-class variance in the change-index distribution, and this
stratum-specific threshold is then constrained by global and
class-specific lower bounds to prevent over-permissive detection in
spectrally noisy strata.

Pixels above the resulting seed threshold are treated as high-confidence
burned seeds and are subsequently expanded into contiguous patches
through region growing under a more permissive class-specific growth
threshold. Patches are retained only when they contain a minimum number
of seed pixels, which ensures that every mapped patch is anchored to a
robust spectral core and limits the influence of isolated false
positives. The result is an initial raster of seed-grown candidate
patches.

Each candidate patch is then evaluated by a sequence of rule-based
filters. The first filter assesses internal support, including seed
evidence, minimum patch plausibility, dominance of burnable land-cover
classes and the absence of excessive missing or non-burnable cover. The
second filter evaluates external corroboration using active-fire
hotspots when reliable hotspot detections are available for the target
year; when they are not, this filter is recorded as not applied, so that
pre-hotspot years are not penalised for the absence of an input that
does not exist for them. The third filter assesses spectral consistency
by comparing each candidate patch against a burned-like reference
distribution derived from high-confidence patches. Together, these
filters integrate into a canonical unsupervised decision layer that
assigns every patch to one of three classes: \textbf{keep}, indicating
high-confidence burned candidates with strong internal or external
support; \textbf{drop}, indicating likely unburned areas or false
positives; and \textbf{review}, indicating ambiguous cases often
associated with haze, shadows, drought-stressed herbaceous systems,
land-water adjacency effects or other persistent spectral confounders.

All filter outcomes, support metrics and stratification metadata are
stored as polygon-level attributes. The full decision path of every
filter is logged at the patch level, so that the reasons behind each
keep, drop or review assignment can be reconstructed afterwards. This is
what makes the rule-based unsupervised phase auditable: the
candidate-patch layer is not only a map but also a record of the spatial
reasoning that produced it.

\subsection{Supervised phase}\label{supervised-phase}

The supervised phase refines the candidate-patch universe produced by
the unsupervised phase by assigning a burned probability to each patch.
Its purpose is to reduce residual omission and commission errors in
spectrally ambiguous settings without redefining the candidate universe
and without depending on externally compiled training labels.

Training pools are derived directly from the unsupervised output.
High-confidence keep patches form the burned training pool, while the
unburned training pool is assembled from complementary negative
evidence: drop patches rejected by the rule-based filters, random
background samples drawn from burnable areas subject to exclusion
buffers around keep patches and previously burned areas, and
lower-confidence Otsu-derived candidates. The contribution of each
source is capped relative to the burned pool to prevent class imbalance
from biasing the classifier. Review patches are withheld from training,
since they represent the cases for which unsupervised evidence is
ambiguous, and are scored only at inference time.

Patch-level predictors are aggregated from pixel-level information
within each candidate patch. They include spectral summaries derived
from the immediate and, where applicable, delayed change-index
composites; contextual descriptors capturing vegetation composition,
ecoregion membership, elevation and slope; and active-fire variables
describing the presence, intensity and spatial configuration of hotspot
detections when these are available. Day-of-year metrics associated with
the selected change-index observations are included as temporal
predictors.

The supervised classifier is implemented as a gradient-boosted tree
model (XGBoost) and trained under spatial block cross-validation to
mitigate spatial autocorrelation between training and validation splits.
The framework supports three complementary modelling variants. The
\textbf{one-year} variant trains and evaluates within the same calendar
year and provides the operational within-year baseline. The
\textbf{out-of-fold} variant uses spatial cross-validation within a
single year to produce predictions for which each patch has been scored
by a model that did not see it during training. The
\textbf{leave-one-year-out} variant trains on the remaining years and
predicts on the held-out year, providing the strict cross-year transfer
estimate.

Because the unsupervised phase has already filtered out most non-burned
candidates, the supervised stage operates on a pre-curated universe in
which the majority of patches are already plausible burned candidates.
Predicted probabilities therefore tend to concentrate close to one, with
the operational decision threshold also lying closer to one than to the
conventional 0.5 cut-off. The supervised stage in OtsuFire is
accordingly best understood as a re-ranker over a pre-curated candidate
set rather than as an independent burned-area classifier.

\section{Inputs}\label{inputs}

OtsuFire derives patch-level information from six conceptual input
groups that jointly support the unsupervised filters and the supervised
predictors. These groups are summarised below; their formal definition,
sources and detailed variable lists are provided in the Supplementary
Material of the OtsuFire paper.

\textbf{Immediate change-index composites} are derived from
medium-resolution optical imagery (Landsat for earlier years and
Sentinel-2 from 2018 onwards), restricted to the summer fire season (1
June to 30 September). For each target year, pre- and post-fire
composites are obtained by selecting the minimum NBR observation within
predefined temporal windows, and RBR is computed from these temporally
matched composites as the primary measure of spectral change. The
framework is index-agnostic; RBR is used as the calibration reference
for the implementation described here.

\textbf{Delayed change-index composites} are derived from autumn-winter
observations following the target fire season. They capture persistence
of the spectral change beyond the immediate post-fire period and help to
separate true burned areas from transient signals such as crop
harvesting or seasonal drying.

\textbf{Land-cover information} is taken from CORINE Land Cover
products, reclassified into eleven ecologically relevant classes. These
classes underpin three workflow components: the burnable mask that
constrains both candidate generation and validation, the vegetation
strata used in combination with ecoregions to stratify Otsu
thresholding, and the categorical vegetation predictors used in the
supervised phase. Ecoregions are derived from the WWF Terrestrial
Ecoregions of the World and are intersected with the vegetation classes
to define ecologically homogeneous units for stratification.

\textbf{Topographic information} comprises elevation and slope derived
from a 30-metre digital elevation model. Topography is used as a set of
continuous predictors in the supervised phase, where it helps account
for systematic variations in illumination, moisture availability and
vegetation structure, but it is not used to stratify the unsupervised
phase.

\textbf{Active-fire information} is obtained from satellite hotspot
detections, available reliably from 1995 onwards. Hotspot variables
describe temporal availability, spatial presence, intensity statistics
and contextual signals at the patch level, and play two complementary
roles in the workflow. In the unsupervised phase, they contribute to the
external corroboration filter, providing additional evidence for
candidate patches intersected by, or located near, retained detections.
In the supervised phase, they enter as patch-level predictors alongside
spectral, contextual and temporal variables. Because hotspot
availability varies across the historical period, a
temporal-availability flag is used to distinguish years with reliable
detections from periods predating the products, allowing the workflow to
remain applicable to the full historical time series.

\textbf{Day-of-year metrics} summarise the acquisition dates of the
observations selected for the immediate and delayed change-index
composites. They are included as temporal predictors in the supervised
phase.

\textbf{Reference burned-area perimeters} from official national,
regional and pan-European sources (EFFIS) are used exclusively for
validation. They are treated as independent external benchmarks rather
than as training labels, and are filtered to the summer fire season and
to burnable land-cover classes to ensure comparability with the mapped
outputs.

\bookmarksetup{startatroot}

\chapter{How to use the OtsuFire
package}\label{how-to-use-the-otsufire-package}

\section{Introduction: installing and loading the
package}\label{introduction-installing-and-loading-the-package}

OtsuFire is distributed through the Comprehensive R Archive Network
(CRAN) and can be installed in the same way as any other R package. The
recommended workflow is to install it once from the official source,
attach it at the start of each session, and then call its functions
directly.

To install the package, run the following inside an R session connected
to the internet:

\begin{Shaded}
\begin{Highlighting}[]
\FunctionTok{install.packages}\NormalTok{(}\StringTok{"OtsuFire"}\NormalTok{)}
\end{Highlighting}
\end{Shaded}

This command pulls the latest released version from CRAN, together with
its mandatory dependencies, and places the package in the default user
library. Once installed, the package is loaded for use with:

\begin{Shaded}
\begin{Highlighting}[]
\FunctionTok{library}\NormalTok{(OtsuFire)}
\end{Highlighting}
\end{Shaded}

Two ancillary R packages are used throughout the tutorial alongside
OtsuFire. The \texttt{sf} package provides vector data handling, while
\texttt{terra} provides raster handling. Both are loaded explicitly at
the top of each worked example:

\begin{Shaded}
\begin{Highlighting}[]
\FunctionTok{library}\NormalTok{(sf)}
\FunctionTok{library}\NormalTok{(terra)}
\end{Highlighting}
\end{Shaded}

Some processing steps in the unsupervised workflow delegate raster and
vector operations to external GDAL utilities (\texttt{gdal\_polygonize},
\texttt{gdalwarp}, \texttt{ogr2ogr}) and to a Python interpreter. These
tools must be available on the host system; OtsuFire is then told where
to find them through configuration paths set later in this tutorial.

Once OtsuFire is installed and loaded, the package is ready to use. The
next sections illustrate its workflow through worked examples, starting
with the unsupervised stage.

\section{Unsupervised stage}\label{unsupervised-stage-1}

The unsupervised stage of OtsuFire converts an annual spectral-change
composite into an auditable layer of candidate burned patches. It is
composed of three operational steps: building a configuration object
that bundles inputs and parameters, detecting candidate patches through
stratified Otsu thresholding and seed-and-grow segmentation, and scoring
those candidates through rule-based filtering into the canonical keep /
review / drop decision layer.

The behaviour of the stage depends on the data available for the target
year. In particular, the second rule-based filter relies on active-fire
hotspot detections, which are only available from 1995 onwards. The
tutorial therefore separates two cases: years with active hotspot
support, where the full rule-based system can be applied, and earlier
years, for which the workflow records the hotspot filter as not applied
rather than as failed. The first case is illustrated below; earlier
years will be covered in a later subsection.

\subsection{Year with active hotspot
support}\label{year-with-active-hotspot-support}

This worked example runs the unsupervised stage for a target year in
which active-fire hotspot detections are available. The example uses
\textbf{2025} as the target year and the \textbf{balanced} detection
preset, but the same script applies to any year from 1995 onwards by
changing the two user settings at the top of the user-settings block.

The example is presented here as a self-contained draft. Subsequent
revisions of this chapter will execute it block by block and discuss the
intermediate outputs in detail.

\subsubsection{Package load}\label{package-load}

\begin{Shaded}
\begin{Highlighting}[]
\FunctionTok{suppressPackageStartupMessages}\NormalTok{(\{}
  \FunctionTok{library}\NormalTok{(OtsuFire)}
  \FunctionTok{library}\NormalTok{(sf)}
  \FunctionTok{library}\NormalTok{(terra)}
\NormalTok{\})}
\end{Highlighting}
\end{Shaded}

\subsubsection{User settings}\label{user-settings}

Choose one target year and one detection preset. Everything below is
derived from these two values.

\begin{Shaded}
\begin{Highlighting}[]
\NormalTok{base\_dir }\OtherTok{\textless{}{-}} \StringTok{"C:/00\_NATALIA\_DOCTORADO/00\_FIRE\_MAPPING"}

\NormalTok{target\_year }\OtherTok{\textless{}{-}} \DecValTok{2025}\NormalTok{L}

\NormalTok{selected\_preset }\OtherTok{\textless{}{-}} \StringTok{"balanced"}

\CommentTok{\# Typical options:}
\CommentTok{\#   "permissive"}
\CommentTok{\#   "balanced"}
\CommentTok{\#   "conservative"}
\CommentTok{\#   "very\_conservative"}
\end{Highlighting}
\end{Shaded}

\subsubsection{Tool paths}\label{tool-paths}

These paths are needed because the unsupervised workflow delegates some
raster and vector operations to external GDAL / Python tools.

\begin{Shaded}
\begin{Highlighting}[]
\NormalTok{python\_exe }\OtherTok{\textless{}{-}} \StringTok{"C:/Users/Olga.Viedma/AppData/Local/Anaconda3/python.exe"}
\NormalTok{gdal\_polygonize\_script }\OtherTok{\textless{}{-}} \StringTok{"C:/Users/Olga.Viedma/AppData/Local/Anaconda3/Scripts/gdal\_polygonize.py"}
\NormalTok{gdalwarp\_path }\OtherTok{\textless{}{-}} \StringTok{"C:/Users/Olga.Viedma/AppData/Local/Anaconda3/Library/bin/gdalwarp.exe"}
\NormalTok{ogr2ogr\_exe }\OtherTok{\textless{}{-}} \StringTok{"C:/Users/Olga.Viedma/AppData/Local/Anaconda3/Library/bin/ogr2ogr.exe"}
\end{Highlighting}
\end{Shaded}

\subsubsection{Derived year settings}\label{derived-year-settings}

CORINE land-cover maps are not available every year. This block maps the
target fire year to the corresponding CORINE reference year used by the
workflow.

Rule:

\begin{itemize}
\tightlist
\item
  \texttt{target\_year\ \textless{}\ 2000} → CORINE 1990
\item
  \texttt{2000–2005} → CORINE 2000
\item
  \texttt{2006–2011} → CORINE 2006
\item
  \texttt{2012–2017} → CORINE 2012
\item
  \texttt{2018\ onwards} → CORINE 2018
\end{itemize}

\begin{Shaded}
\begin{Highlighting}[]
\NormalTok{corine\_year }\OtherTok{\textless{}{-}} \ControlFlowTok{if}\NormalTok{ (target\_year }\SpecialCharTok{\textless{}} \DecValTok{2000}\NormalTok{) \{}
  \DecValTok{1990}
\NormalTok{\} }\ControlFlowTok{else} \ControlFlowTok{if}\NormalTok{ (target\_year }\SpecialCharTok{\textless{}} \DecValTok{2006}\NormalTok{) \{}
  \DecValTok{2000}
\NormalTok{\} }\ControlFlowTok{else} \ControlFlowTok{if}\NormalTok{ (target\_year }\SpecialCharTok{\textless{}} \DecValTok{2012}\NormalTok{) \{}
  \DecValTok{2006}
\NormalTok{\} }\ControlFlowTok{else} \ControlFlowTok{if}\NormalTok{ (target\_year }\SpecialCharTok{\textless{}} \DecValTok{2018}\NormalTok{) \{}
  \DecValTok{2012}
\NormalTok{\} }\ControlFlowTok{else}\NormalTok{ \{}
  \DecValTok{2018}
\NormalTok{\}}
\end{Highlighting}
\end{Shaded}

\subsubsection{Input paths}\label{input-paths}

These are the spatial inputs consumed by the unsupervised workflow.

Some paths depend on \texttt{target\_year} (\texttt{change\_index},
\texttt{reference\_burned\_map}, \texttt{previous\_year\_burned},
\texttt{hotspots}). Others depend on \texttt{corine\_year}
(\texttt{vegetation\_map}, \texttt{burnable\_mask}).

\begin{Shaded}
\begin{Highlighting}[]
\NormalTok{change\_index }\OtherTok{\textless{}{-}} \FunctionTok{file.path}\NormalTok{(}
\NormalTok{  base\_dir,}
  \StringTok{"1\_DATA/Imagery/Composites\_90m/Min\_Min"}\NormalTok{,}
  \FunctionTok{paste0}\NormalTok{(}\StringTok{"MinMin\_"}\NormalTok{, target\_year, }\StringTok{"\_mosaic\_res90m.tif"}\NormalTok{)}
\NormalTok{)}

\NormalTok{vegetation\_map }\OtherTok{\textless{}{-}} \FunctionTok{file.path}\NormalTok{(}
\NormalTok{  base\_dir,}
  \StringTok{"1\_DATA/Corine\_Masks"}\NormalTok{,}
  \FunctionTok{paste0}\NormalTok{(}\StringTok{"CLC\_"}\NormalTok{, corine\_year, }\StringTok{"\_peninsula.tif"}\NormalTok{)}
\NormalTok{)}

\NormalTok{burnable\_mask }\OtherTok{\textless{}{-}} \FunctionTok{file.path}\NormalTok{(}
\NormalTok{  base\_dir,}
  \StringTok{"1\_DATA/Corine\_Masks"}\NormalTok{,}
  \FunctionTok{paste0}\NormalTok{(}
    \StringTok{"burneable\_mask\_binary\_corine"}\NormalTok{,}
    \FunctionTok{substr}\NormalTok{(corine\_year, }\DecValTok{3}\NormalTok{, }\DecValTok{4}\NormalTok{),}
    \StringTok{"\_wgs84.tif"}
\NormalTok{  )}
\NormalTok{)}

\NormalTok{reference\_burned\_map }\OtherTok{\textless{}{-}} \FunctionTok{file.path}\NormalTok{(}
\NormalTok{  base\_dir,}
  \StringTok{"1\_DATA/Fires/Validation\_fires\_burneable\_verano"}\NormalTok{,}
  \FunctionTok{paste0}\NormalTok{(}\StringTok{"Effis\_CA\_"}\NormalTok{, target\_year, }\StringTok{"\_maskKeep\_summer.shp"}\NormalTok{)}
\NormalTok{)}

\NormalTok{previous\_year\_burned }\OtherTok{\textless{}{-}} \ControlFlowTok{if}\NormalTok{ (target\_year }\SpecialCharTok{\textgreater{}} \DecValTok{1985}\NormalTok{) \{}
  \FunctionTok{file.path}\NormalTok{(}
\NormalTok{    base\_dir,}
    \StringTok{"1\_DATA/Fires/Validation\_fires"}\NormalTok{,}
    \FunctionTok{paste0}\NormalTok{(}\StringTok{"Effis\_CA\_"}\NormalTok{, target\_year }\SpecialCharTok{{-}} \DecValTok{1}\NormalTok{, }\StringTok{".shp"}\NormalTok{)}
\NormalTok{  )}
\NormalTok{\} }\ControlFlowTok{else}\NormalTok{ \{}
  \ConstantTok{NULL}
\NormalTok{\}}

\NormalTok{hotspots }\OtherTok{\textless{}{-}} \ControlFlowTok{if}\NormalTok{ (target\_year }\SpecialCharTok{\textless{}} \DecValTok{1995}\NormalTok{) \{}
  \ConstantTok{NULL}
\NormalTok{\} }\ControlFlowTok{else}\NormalTok{ \{}
  \FunctionTok{file.path}\NormalTok{(}
\NormalTok{    base\_dir,}
    \StringTok{"1\_DATA/Hotspots"}\NormalTok{,}
    \FunctionTok{paste0}\NormalTok{(}\StringTok{"hotspots\_iberia\_"}\NormalTok{, target\_year, }\StringTok{".geojson"}\NormalTok{)}
\NormalTok{  )}
\NormalTok{\}}

\NormalTok{ecoregion\_shapefile\_path }\OtherTok{\textless{}{-}} \FunctionTok{file.path}\NormalTok{(}
\NormalTok{  base\_dir,}
  \StringTok{"1\_DATA/Ecoregion/ecoregiones\_olson.shp"}
\NormalTok{)}

\NormalTok{peninsula\_shapefile\_path }\OtherTok{\textless{}{-}} \FunctionTok{file.path}\NormalTok{(}
\NormalTok{  base\_dir,}
  \StringTok{"1\_DATA/Borders/Iberian\_Peninsula.shp"}
\NormalTok{)}
\end{Highlighting}
\end{Shaded}

\subsubsection{Quick input check}\label{quick-input-check}

This check verifies that the expected files exist before the workflow
starts. Optional inputs (\texttt{previous\_year\_burned} may be NULL for
1985, \texttt{hotspots} is NULL for pre-hotspot years) are excluded from
the failure condition.

\begin{Shaded}
\begin{Highlighting}[]
\NormalTok{paths\_to\_check }\OtherTok{\textless{}{-}} \FunctionTok{list}\NormalTok{(}
  \AttributeTok{change\_index =}\NormalTok{ change\_index,}
  \AttributeTok{vegetation\_map =}\NormalTok{ vegetation\_map,}
  \AttributeTok{burnable\_mask =}\NormalTok{ burnable\_mask,}
  \AttributeTok{reference\_burned\_map =}\NormalTok{ reference\_burned\_map,}
  \AttributeTok{previous\_year\_burned =}\NormalTok{ previous\_year\_burned,}
  \AttributeTok{hotspots =}\NormalTok{ hotspots,}
  \AttributeTok{ecoregion\_shapefile\_path =}\NormalTok{ ecoregion\_shapefile\_path,}
  \AttributeTok{peninsula\_shapefile\_path =}\NormalTok{ peninsula\_shapefile\_path}
\NormalTok{)}

\NormalTok{input\_check }\OtherTok{\textless{}{-}} \FunctionTok{data.frame}\NormalTok{(}
  \AttributeTok{input =} \FunctionTok{names}\NormalTok{(paths\_to\_check),}
  \AttributeTok{path =} \FunctionTok{vapply}\NormalTok{(}
\NormalTok{    paths\_to\_check,}
    \ControlFlowTok{function}\NormalTok{(x) }\ControlFlowTok{if}\NormalTok{ (}\FunctionTok{is.null}\NormalTok{(x)) }\ConstantTok{NA\_character\_} \ControlFlowTok{else}\NormalTok{ x,}
    \FunctionTok{character}\NormalTok{(}\DecValTok{1}\NormalTok{)}
\NormalTok{  ),}
  \AttributeTok{exists =} \FunctionTok{vapply}\NormalTok{(}
\NormalTok{    paths\_to\_check,}
    \ControlFlowTok{function}\NormalTok{(x) }\ControlFlowTok{if}\NormalTok{ (}\FunctionTok{is.null}\NormalTok{(x)) }\ConstantTok{NA} \ControlFlowTok{else} \FunctionTok{file.exists}\NormalTok{(x),}
    \FunctionTok{logical}\NormalTok{(}\DecValTok{1}\NormalTok{)}
\NormalTok{  )}
\NormalTok{)}

\NormalTok{input\_check}

\NormalTok{required\_inputs }\OtherTok{\textless{}{-}} \FunctionTok{c}\NormalTok{(}
  \StringTok{"change\_index"}\NormalTok{,}
  \StringTok{"vegetation\_map"}\NormalTok{,}
  \StringTok{"burnable\_mask"}\NormalTok{,}
  \StringTok{"reference\_burned\_map"}\NormalTok{,}
  \StringTok{"ecoregion\_shapefile\_path"}\NormalTok{,}
  \StringTok{"peninsula\_shapefile\_path"}
\NormalTok{)}

\ControlFlowTok{if}\NormalTok{ (}\FunctionTok{any}\NormalTok{(}\SpecialCharTok{!}\NormalTok{input\_check}\SpecialCharTok{$}\NormalTok{exists[input\_check}\SpecialCharTok{$}\NormalTok{input }\SpecialCharTok{\%in\%}\NormalTok{ required\_inputs])) \{}
  \FunctionTok{stop}\NormalTok{(}\StringTok{"Some required input files are missing. Check input\_check."}\NormalTok{)}
\NormalTok{\}}
\end{Highlighting}
\end{Shaded}

\subsubsection{Detection presets}\label{detection-presets}

These are the main detection knobs.

\begin{itemize}
\tightlist
\item
  \texttt{seed\_threshold} --- minimum change value needed to create a
  burned seed pixel. Higher values mean fewer seeds and more
  conservative detection.
\item
  \texttt{growth\_delta} --- amount by which the grow threshold is
  relaxed relative to the seed threshold. Higher values mean easier
  patch expansion.
\item
  \texttt{minimum\_growth\_threshold} --- absolute floor below which
  region growing is not allowed to go. Higher values mean tighter, more
  conservative patch growth.
\item
  \texttt{minimum\_seed\_pixels} --- minimum number of seed pixels
  needed to keep a candidate component. Higher values remove small
  isolated detections.
\item
  \texttt{*\_by\_vegetation} --- per-class overrides for the global
  values above, used when one global value is too coarse and vegetation
  types need different behaviour.
\end{itemize}

\begin{Shaded}
\begin{Highlighting}[]
\NormalTok{preset\_params }\OtherTok{\textless{}{-}} \FunctionTok{list}\NormalTok{(}

  \AttributeTok{permissive =} \FunctionTok{list}\NormalTok{(}
    \AttributeTok{expected =} \StringTok{"Highest sensitivity. More candidate burned patches, but higher commission risk."}\NormalTok{,}
    \AttributeTok{detect\_params =} \FunctionTok{list}\NormalTok{(}
      \AttributeTok{seed\_threshold =} \DecValTok{300}\NormalTok{,}
      \AttributeTok{growth\_delta =} \DecValTok{100}\NormalTok{,}
      \AttributeTok{minimum\_growth\_threshold =} \DecValTok{200}\NormalTok{,}
      \AttributeTok{minimum\_seed\_pixels =} \DecValTok{20}\NormalTok{L,}
      \AttributeTok{seed\_threshold\_by\_vegetation =} \FunctionTok{c}\NormalTok{(}
        \StringTok{"1"} \OtherTok{=} \DecValTok{300}\NormalTok{, }\StringTok{"2"} \OtherTok{=} \DecValTok{350}\NormalTok{, }\StringTok{"3"} \OtherTok{=} \DecValTok{350}\NormalTok{, }\StringTok{"4"} \OtherTok{=} \DecValTok{350}\NormalTok{, }\StringTok{"5"} \OtherTok{=} \DecValTok{350}\NormalTok{,}
        \StringTok{"6"} \OtherTok{=} \DecValTok{350}\NormalTok{, }\StringTok{"7"} \OtherTok{=} \DecValTok{350}\NormalTok{, }\StringTok{"8"} \OtherTok{=} \DecValTok{500}\NormalTok{, }\StringTok{"9"} \OtherTok{=} \DecValTok{440}\NormalTok{, }\StringTok{"10"} \OtherTok{=} \DecValTok{650}\NormalTok{, }\StringTok{"11"} \OtherTok{=} \DecValTok{450}
\NormalTok{      ),}
      \AttributeTok{growth\_delta\_by\_vegetation =} \FunctionTok{c}\NormalTok{(}
        \StringTok{"1"} \OtherTok{=} \DecValTok{100}\NormalTok{, }\StringTok{"2"} \OtherTok{=} \DecValTok{50}\NormalTok{, }\StringTok{"3"} \OtherTok{=} \DecValTok{100}\NormalTok{, }\StringTok{"4"} \OtherTok{=} \DecValTok{120}\NormalTok{, }\StringTok{"5"} \OtherTok{=} \DecValTok{100}\NormalTok{,}
        \StringTok{"6"} \OtherTok{=} \DecValTok{100}\NormalTok{, }\StringTok{"7"} \OtherTok{=} \DecValTok{100}\NormalTok{, }\StringTok{"8"} \OtherTok{=} \DecValTok{10}\NormalTok{, }\StringTok{"9"} \OtherTok{=} \DecValTok{15}\NormalTok{, }\StringTok{"10"} \OtherTok{=} \DecValTok{5}\NormalTok{, }\StringTok{"11"} \OtherTok{=} \DecValTok{15}
\NormalTok{      ),}
      \AttributeTok{minimum\_growth\_threshold\_by\_vegetation =} \FunctionTok{c}\NormalTok{(}
        \StringTok{"1"} \OtherTok{=} \DecValTok{200}\NormalTok{, }\StringTok{"2"} \OtherTok{=} \DecValTok{300}\NormalTok{, }\StringTok{"3"} \OtherTok{=} \DecValTok{250}\NormalTok{, }\StringTok{"4"} \OtherTok{=} \DecValTok{210}\NormalTok{, }\StringTok{"5"} \OtherTok{=} \DecValTok{220}\NormalTok{,}
        \StringTok{"6"} \OtherTok{=} \DecValTok{220}\NormalTok{, }\StringTok{"7"} \OtherTok{=} \DecValTok{220}\NormalTok{, }\StringTok{"8"} \OtherTok{=} \DecValTok{400}\NormalTok{, }\StringTok{"9"} \OtherTok{=} \DecValTok{380}\NormalTok{, }\StringTok{"10"} \OtherTok{=} \DecValTok{550}\NormalTok{, }\StringTok{"11"} \OtherTok{=} \DecValTok{380}
\NormalTok{      )}
\NormalTok{    )}
\NormalTok{  ),}

  \AttributeTok{balanced =} \FunctionTok{list}\NormalTok{(}
    \AttributeTok{expected =} \StringTok{"Default operational compromise between omission and commission."}\NormalTok{,}
    \AttributeTok{detect\_params =} \FunctionTok{list}\NormalTok{(}
      \AttributeTok{seed\_threshold =} \DecValTok{310}\NormalTok{,}
      \AttributeTok{growth\_delta =} \DecValTok{90}\NormalTok{,}
      \AttributeTok{minimum\_growth\_threshold =} \DecValTok{240}\NormalTok{,}
      \AttributeTok{minimum\_seed\_pixels =} \DecValTok{30}\NormalTok{L,}
      \AttributeTok{seed\_threshold\_by\_vegetation =} \FunctionTok{c}\NormalTok{(}
        \StringTok{"1"} \OtherTok{=} \DecValTok{350}\NormalTok{, }\StringTok{"2"} \OtherTok{=} \DecValTok{420}\NormalTok{, }\StringTok{"3"} \OtherTok{=} \DecValTok{380}\NormalTok{, }\StringTok{"4"} \OtherTok{=} \DecValTok{320}\NormalTok{, }\StringTok{"5"} \OtherTok{=} \DecValTok{320}\NormalTok{,}
        \StringTok{"6"} \OtherTok{=} \DecValTok{320}\NormalTok{, }\StringTok{"7"} \OtherTok{=} \DecValTok{320}\NormalTok{, }\StringTok{"8"} \OtherTok{=} \DecValTok{500}\NormalTok{, }\StringTok{"9"} \OtherTok{=} \DecValTok{470}\NormalTok{, }\StringTok{"10"} \OtherTok{=} \DecValTok{670}\NormalTok{, }\StringTok{"11"} \OtherTok{=} \DecValTok{480}
\NormalTok{      ),}
      \AttributeTok{growth\_delta\_by\_vegetation =} \FunctionTok{c}\NormalTok{(}
        \StringTok{"1"} \OtherTok{=} \DecValTok{70}\NormalTok{, }\StringTok{"2"} \OtherTok{=} \DecValTok{20}\NormalTok{, }\StringTok{"3"} \OtherTok{=} \DecValTok{40}\NormalTok{, }\StringTok{"4"} \OtherTok{=} \DecValTok{100}\NormalTok{, }\StringTok{"5"} \OtherTok{=} \DecValTok{85}\NormalTok{,}
        \StringTok{"6"} \OtherTok{=} \DecValTok{85}\NormalTok{, }\StringTok{"7"} \OtherTok{=} \DecValTok{85}\NormalTok{, }\StringTok{"8"} \OtherTok{=} \DecValTok{10}\NormalTok{, }\StringTok{"9"} \OtherTok{=} \DecValTok{10}\NormalTok{, }\StringTok{"10"} \OtherTok{=} \DecValTok{5}\NormalTok{, }\StringTok{"11"} \OtherTok{=} \DecValTok{10}
\NormalTok{      ),}
      \AttributeTok{minimum\_growth\_threshold\_by\_vegetation =} \FunctionTok{c}\NormalTok{(}
        \StringTok{"1"} \OtherTok{=} \DecValTok{280}\NormalTok{, }\StringTok{"2"} \OtherTok{=} \DecValTok{360}\NormalTok{, }\StringTok{"3"} \OtherTok{=} \DecValTok{300}\NormalTok{, }\StringTok{"4"} \OtherTok{=} \DecValTok{230}\NormalTok{, }\StringTok{"5"} \OtherTok{=} \DecValTok{240}\NormalTok{,}
        \StringTok{"6"} \OtherTok{=} \DecValTok{240}\NormalTok{, }\StringTok{"7"} \OtherTok{=} \DecValTok{240}\NormalTok{, }\StringTok{"8"} \OtherTok{=} \DecValTok{430}\NormalTok{, }\StringTok{"9"} \OtherTok{=} \DecValTok{400}\NormalTok{, }\StringTok{"10"} \OtherTok{=} \DecValTok{560}\NormalTok{, }\StringTok{"11"} \OtherTok{=} \DecValTok{400}
\NormalTok{      )}
\NormalTok{    )}
\NormalTok{  ),}

  \AttributeTok{conservative =} \FunctionTok{list}\NormalTok{(}
    \AttributeTok{expected =} \StringTok{"More restrictive. Reduces false positives and over{-}expansion, with increased omission risk."}\NormalTok{,}
    \AttributeTok{detect\_params =} \FunctionTok{list}\NormalTok{(}
      \AttributeTok{seed\_threshold =} \DecValTok{400}\NormalTok{,}
      \AttributeTok{growth\_delta =} \DecValTok{50}\NormalTok{,}
      \AttributeTok{minimum\_growth\_threshold =} \DecValTok{320}\NormalTok{,}
      \AttributeTok{minimum\_seed\_pixels =} \DecValTok{50}\NormalTok{L,}
      \AttributeTok{seed\_threshold\_by\_vegetation =} \FunctionTok{c}\NormalTok{(}
        \StringTok{"1"} \OtherTok{=} \DecValTok{400}\NormalTok{, }\StringTok{"2"} \OtherTok{=} \DecValTok{470}\NormalTok{, }\StringTok{"3"} \OtherTok{=} \DecValTok{430}\NormalTok{, }\StringTok{"4"} \OtherTok{=} \DecValTok{400}\NormalTok{, }\StringTok{"5"} \OtherTok{=} \DecValTok{400}\NormalTok{,}
        \StringTok{"6"} \OtherTok{=} \DecValTok{400}\NormalTok{, }\StringTok{"7"} \OtherTok{=} \DecValTok{400}\NormalTok{, }\StringTok{"8"} \OtherTok{=} \DecValTok{520}\NormalTok{, }\StringTok{"9"} \OtherTok{=} \DecValTok{500}\NormalTok{, }\StringTok{"10"} \OtherTok{=} \DecValTok{700}\NormalTok{, }\StringTok{"11"} \OtherTok{=} \DecValTok{520}
\NormalTok{      ),}
      \AttributeTok{growth\_delta\_by\_vegetation =} \FunctionTok{c}\NormalTok{(}
        \StringTok{"1"} \OtherTok{=} \DecValTok{50}\NormalTok{, }\StringTok{"2"} \OtherTok{=} \DecValTok{10}\NormalTok{, }\StringTok{"3"} \OtherTok{=} \DecValTok{20}\NormalTok{, }\StringTok{"4"} \OtherTok{=} \DecValTok{60}\NormalTok{, }\StringTok{"5"} \OtherTok{=} \DecValTok{50}\NormalTok{,}
        \StringTok{"6"} \OtherTok{=} \DecValTok{50}\NormalTok{, }\StringTok{"7"} \OtherTok{=} \DecValTok{50}\NormalTok{, }\StringTok{"8"} \OtherTok{=} \DecValTok{5}\NormalTok{, }\StringTok{"9"} \OtherTok{=} \DecValTok{5}\NormalTok{, }\StringTok{"10"} \OtherTok{=} \DecValTok{5}\NormalTok{, }\StringTok{"11"} \OtherTok{=} \DecValTok{5}
\NormalTok{      ),}
      \AttributeTok{minimum\_growth\_threshold\_by\_vegetation =} \FunctionTok{c}\NormalTok{(}
        \StringTok{"1"} \OtherTok{=} \DecValTok{330}\NormalTok{, }\StringTok{"2"} \OtherTok{=} \DecValTok{400}\NormalTok{, }\StringTok{"3"} \OtherTok{=} \DecValTok{370}\NormalTok{, }\StringTok{"4"} \OtherTok{=} \DecValTok{330}\NormalTok{, }\StringTok{"5"} \OtherTok{=} \DecValTok{330}\NormalTok{,}
        \StringTok{"6"} \OtherTok{=} \DecValTok{330}\NormalTok{, }\StringTok{"7"} \OtherTok{=} \DecValTok{330}\NormalTok{, }\StringTok{"8"} \OtherTok{=} \DecValTok{450}\NormalTok{, }\StringTok{"9"} \OtherTok{=} \DecValTok{430}\NormalTok{, }\StringTok{"10"} \OtherTok{=} \DecValTok{580}\NormalTok{, }\StringTok{"11"} \OtherTok{=} \DecValTok{430}
\NormalTok{      )}
\NormalTok{    )}
\NormalTok{  ),}

  \AttributeTok{very\_conservative =} \FunctionTok{list}\NormalTok{(}
    \AttributeTok{expected =} \FunctionTok{paste}\NormalTok{(}
      \StringTok{"Intermediate anti{-}banding preset for problematic historical years."}\NormalTok{,}
      \StringTok{"Designed to suppress Landsat banding artefacts while preserving"}\NormalTok{,}
      \StringTok{"medium and large real burned patches."}
\NormalTok{    ),}
    \AttributeTok{detect\_params =} \FunctionTok{list}\NormalTok{(}
      \AttributeTok{seed\_threshold =} \DecValTok{430}\NormalTok{,}
      \AttributeTok{growth\_delta =} \DecValTok{40}\NormalTok{,}
      \AttributeTok{minimum\_growth\_threshold =} \DecValTok{350}\NormalTok{,}
      \AttributeTok{minimum\_seed\_pixels =} \DecValTok{70}\NormalTok{L,}
      \AttributeTok{seed\_threshold\_by\_vegetation =} \FunctionTok{c}\NormalTok{(}
        \StringTok{"1"} \OtherTok{=} \DecValTok{430}\NormalTok{, }\StringTok{"2"} \OtherTok{=} \DecValTok{480}\NormalTok{, }\StringTok{"3"} \OtherTok{=} \DecValTok{450}\NormalTok{, }\StringTok{"4"} \OtherTok{=} \DecValTok{430}\NormalTok{, }\StringTok{"5"} \OtherTok{=} \DecValTok{430}\NormalTok{,}
        \StringTok{"6"} \OtherTok{=} \DecValTok{430}\NormalTok{, }\StringTok{"7"} \OtherTok{=} \DecValTok{430}\NormalTok{, }\StringTok{"8"} \OtherTok{=} \DecValTok{540}\NormalTok{, }\StringTok{"9"} \OtherTok{=} \DecValTok{520}\NormalTok{, }\StringTok{"10"} \OtherTok{=} \DecValTok{720}\NormalTok{, }\StringTok{"11"} \OtherTok{=} \DecValTok{540}
\NormalTok{      ),}
      \AttributeTok{growth\_delta\_by\_vegetation =} \FunctionTok{c}\NormalTok{(}
        \StringTok{"1"} \OtherTok{=} \DecValTok{40}\NormalTok{, }\StringTok{"2"} \OtherTok{=} \DecValTok{10}\NormalTok{, }\StringTok{"3"} \OtherTok{=} \DecValTok{20}\NormalTok{, }\StringTok{"4"} \OtherTok{=} \DecValTok{45}\NormalTok{, }\StringTok{"5"} \OtherTok{=} \DecValTok{40}\NormalTok{,}
        \StringTok{"6"} \OtherTok{=} \DecValTok{40}\NormalTok{, }\StringTok{"7"} \OtherTok{=} \DecValTok{40}\NormalTok{, }\StringTok{"8"} \OtherTok{=} \DecValTok{5}\NormalTok{, }\StringTok{"9"} \OtherTok{=} \DecValTok{5}\NormalTok{, }\StringTok{"10"} \OtherTok{=} \DecValTok{5}\NormalTok{, }\StringTok{"11"} \OtherTok{=} \DecValTok{5}
\NormalTok{      ),}
      \AttributeTok{minimum\_growth\_threshold\_by\_vegetation =} \FunctionTok{c}\NormalTok{(}
        \StringTok{"1"} \OtherTok{=} \DecValTok{360}\NormalTok{, }\StringTok{"2"} \OtherTok{=} \DecValTok{400}\NormalTok{, }\StringTok{"3"} \OtherTok{=} \DecValTok{370}\NormalTok{, }\StringTok{"4"} \OtherTok{=} \DecValTok{360}\NormalTok{, }\StringTok{"5"} \OtherTok{=} \DecValTok{360}\NormalTok{,}
        \StringTok{"6"} \OtherTok{=} \DecValTok{360}\NormalTok{, }\StringTok{"7"} \OtherTok{=} \DecValTok{360}\NormalTok{, }\StringTok{"8"} \OtherTok{=} \DecValTok{470}\NormalTok{, }\StringTok{"9"} \OtherTok{=} \DecValTok{450}\NormalTok{, }\StringTok{"10"} \OtherTok{=} \DecValTok{600}\NormalTok{, }\StringTok{"11"} \OtherTok{=} \DecValTok{450}
\NormalTok{      )}
\NormalTok{    )}
\NormalTok{  )}
\NormalTok{)}

\ControlFlowTok{if}\NormalTok{ (}\SpecialCharTok{!}\NormalTok{selected\_preset }\SpecialCharTok{\%in\%} \FunctionTok{names}\NormalTok{(preset\_params)) \{}
  \FunctionTok{stop}\NormalTok{(}
    \StringTok{"selected\_preset must be one of: "}\NormalTok{,}
    \FunctionTok{paste}\NormalTok{(}\FunctionTok{names}\NormalTok{(preset\_params), }\AttributeTok{collapse =} \StringTok{", "}\NormalTok{)}
\NormalTok{  )}
\NormalTok{\}}

\NormalTok{preset\_params[[selected\_preset]]}\SpecialCharTok{$}\NormalTok{expected}
\NormalTok{preset\_params[[selected\_preset]]}\SpecialCharTok{$}\NormalTok{detect\_params}
\end{Highlighting}
\end{Shaded}

\subsubsection{Refinement parameters}\label{refinement-parameters}

These are the current package defaults written explicitly so the user
can see them and modify them if needed.

\begin{itemize}
\tightlist
\item
  \texttt{aoi\_buffer\_m} --- buffer used around AOI boundaries during
  refinement to reduce edge artefacts.
\item
  \texttt{omission\_buffer\_m} --- reserved cleanup buffer for
  omission-side adjustments. Zero means no extra omission buffer.
\item
  \texttt{minimum\_detected\_area\_m2} --- minimum polygon area kept
  after refinement. Very permissive by default.
\item
  \texttt{merge\_overlaps} --- if TRUE, overlapping refined polygons are
  merged.
\item
  \texttt{merge\_overlaps\_buffer\_m} --- optional buffer before overlap
  merging. Zero means merge only direct overlaps.
\end{itemize}

\begin{Shaded}
\begin{Highlighting}[]
\NormalTok{refine\_params }\OtherTok{\textless{}{-}} \FunctionTok{list}\NormalTok{(}
  \AttributeTok{aoi\_buffer\_m =} \DecValTok{5000}\NormalTok{,}
  \AttributeTok{omission\_buffer\_m =} \DecValTok{0}\NormalTok{,}
  \AttributeTok{minimum\_detected\_area\_m2 =} \DecValTok{10}\NormalTok{,}
  \AttributeTok{merge\_overlaps =} \ConstantTok{TRUE}\NormalTok{,}
  \AttributeTok{merge\_overlaps\_buffer\_m =} \DecValTok{0}
\NormalTok{)}
\end{Highlighting}
\end{Shaded}

\subsubsection{Scoring parameters}\label{scoring-parameters}

These parameters belong to the deterministic scoring stage itself.

\begin{itemize}
\tightlist
\item
  \texttt{support\_buffer\_m} --- buffer used when checking local
  contextual support around candidate patches.
\item
  \texttt{previous\_fire\_exclusion\_buffer\_m} --- buffer around
  previous-year burns used to detect temporal conflicts.
\item
  \texttt{previous\_fire\_cleanup\_buffer\_m} --- cleanup buffer used
  after overlap removal with previous-year burns.
\item
  \texttt{minimum\_remaining\_area\_m2} --- minimum area a patch must
  retain after temporal cleanup.
\item
  \texttt{reference\_buffer\_m} --- buffer used when comparing against
  support-reference areas in scoring.
\end{itemize}

\begin{Shaded}
\begin{Highlighting}[]
\NormalTok{scoring\_params }\OtherTok{\textless{}{-}} \FunctionTok{list}\NormalTok{(}
  \AttributeTok{support\_buffer\_m =} \DecValTok{0}\NormalTok{,}
  \AttributeTok{previous\_fire\_exclusion\_buffer\_m =} \DecValTok{90}\NormalTok{,}
  \AttributeTok{previous\_fire\_cleanup\_buffer\_m =} \DecValTok{90}\NormalTok{,}
  \AttributeTok{minimum\_remaining\_area\_m2 =} \DecValTok{20000}\NormalTok{,}
  \AttributeTok{reference\_buffer\_m =} \DecValTok{90}
\NormalTok{)}
\end{Highlighting}
\end{Shaded}

\subsubsection{Output paths}\label{output-paths}

\texttt{output\_dir} must be the root results folder only. The package
itself appends
\texttt{\textless{}target\_year\textgreater{}/DETERMINISTIC/\textless{}run\_name\textgreater{}},
so neither \texttt{target\_year} nor \texttt{DETERMINISTIC} should be
added manually.

\begin{Shaded}
\begin{Highlighting}[]
\NormalTok{output\_dir }\OtherTok{\textless{}{-}} \FunctionTok{file.path}\NormalTok{(}
\NormalTok{  base\_dir,}
  \StringTok{"1\_DATA/Results"}
\NormalTok{)}

\NormalTok{run\_name }\OtherTok{\textless{}{-}} \FunctionTok{paste0}\NormalTok{(target\_year, }\StringTok{"\_"}\NormalTok{, selected\_preset)}
\end{Highlighting}
\end{Shaded}

\subsubsection{Step 1. Build the unsupervised
configuration}\label{step-1.-build-the-unsupervised-configuration}

At this stage the workflow is only being prepared. Nothing is detected
or scored yet. This is the right moment to inspect the selected preset,
the parameter groups and the input paths.

\begin{Shaded}
\begin{Highlighting}[]
\NormalTok{cfg }\OtherTok{\textless{}{-}} \FunctionTok{build\_burned\_mapping\_config}\NormalTok{(}
  \AttributeTok{change\_index =}\NormalTok{ change\_index,}
  \AttributeTok{vegetation\_map =}\NormalTok{ vegetation\_map,}
  \AttributeTok{burnable\_mask =}\NormalTok{ burnable\_mask,}
  \AttributeTok{hotspots =}\NormalTok{ hotspots,}
  \AttributeTok{previous\_year\_burned =}\NormalTok{ previous\_year\_burned,}
  \AttributeTok{reference\_burned\_map =}\NormalTok{ reference\_burned\_map,}
  \AttributeTok{target\_year =}\NormalTok{ target\_year,}
  \AttributeTok{output\_dir =}\NormalTok{ output\_dir,}
  \AttributeTok{run\_name =}\NormalTok{ run\_name,}
  \AttributeTok{detect\_params =}\NormalTok{ preset\_params[[selected\_preset]]}\SpecialCharTok{$}\NormalTok{detect\_params,}
  \AttributeTok{refine\_params =}\NormalTok{ refine\_params,}
  \AttributeTok{scoring\_params =}\NormalTok{ scoring\_params,}
  \AttributeTok{options =} \FunctionTok{list}\NormalTok{(}
    \AttributeTok{ecoregion\_shapefile\_path =}\NormalTok{ ecoregion\_shapefile\_path,}
    \AttributeTok{peninsula\_shapefile\_path =}\NormalTok{ peninsula\_shapefile\_path,}
    \AttributeTok{change\_index\_validation =} \FunctionTok{list}\NormalTok{(}
      \AttributeTok{expected\_nodata =} \SpecialCharTok{{-}}\DecValTok{9999}\NormalTok{,}
      \AttributeTok{lower\_cap =} \SpecialCharTok{{-}}\DecValTok{1000}\NormalTok{,}
      \AttributeTok{cap\_below =} \ConstantTok{TRUE}\NormalTok{,}
      \AttributeTok{check\_finite =} \ConstantTok{TRUE}
\NormalTok{    )}
\NormalTok{  )}
\NormalTok{)}

\NormalTok{cfg}\SpecialCharTok{$}\NormalTok{tool\_paths}\SpecialCharTok{$}\NormalTok{python\_exe }\OtherTok{\textless{}{-}}\NormalTok{ python\_exe}
\NormalTok{cfg}\SpecialCharTok{$}\NormalTok{tool\_paths}\SpecialCharTok{$}\NormalTok{gdal\_polygonize\_script }\OtherTok{\textless{}{-}}\NormalTok{ gdal\_polygonize\_script}
\NormalTok{cfg}\SpecialCharTok{$}\NormalTok{tool\_paths}\SpecialCharTok{$}\NormalTok{gdalwarp\_path }\OtherTok{\textless{}{-}}\NormalTok{ gdalwarp\_path}
\NormalTok{cfg}\SpecialCharTok{$}\NormalTok{tool\_paths}\SpecialCharTok{$}\NormalTok{ogr2ogr\_exe }\OtherTok{\textless{}{-}}\NormalTok{ ogr2ogr\_exe}
\end{Highlighting}
\end{Shaded}

\subsubsection{Step 2. Detect candidate burned
patches}\label{step-2.-detect-candidate-burned-patches}

Detection uses the chosen preset. The preset controls how easily burned
seed pixels are created, how far candidate patches can grow, and how
much behaviour differs by vegetation / land-cover class.

\begin{Shaded}
\begin{Highlighting}[]
\NormalTok{det }\OtherTok{\textless{}{-}} \FunctionTok{detect\_burned\_patches}\NormalTok{(}
  \AttributeTok{config =}\NormalTok{ cfg,}
  \AttributeTok{write\_outputs =} \ConstantTok{TRUE}\NormalTok{,}
  \AttributeTok{overwrite =} \ConstantTok{TRUE}
\NormalTok{)}

\NormalTok{grown\_candidates\_path }\OtherTok{\textless{}{-}}\NormalTok{ det}\SpecialCharTok{$}\NormalTok{grown\_patches\_path}
\NormalTok{refined\_candidates\_path }\OtherTok{\textless{}{-}}\NormalTok{ det}\SpecialCharTok{$}\NormalTok{refined\_patches\_path}

\FunctionTok{stopifnot}\NormalTok{(}\FunctionTok{file.exists}\NormalTok{(grown\_candidates\_path))}
\FunctionTok{stopifnot}\NormalTok{(}\FunctionTok{file.exists}\NormalTok{(refined\_candidates\_path))}
\end{Highlighting}
\end{Shaded}

\subsubsection{Step 3. Same-year
scoring}\label{step-3.-same-year-scoring}

This is the native deterministic scoring behaviour. Setting
\texttt{keep\_pool\ =\ NULL} tells the package to use its built-in
same-year local fallback.

\begin{Shaded}
\begin{Highlighting}[]
\NormalTok{scored }\OtherTok{\textless{}{-}} \FunctionTok{score\_burned\_patches}\NormalTok{(}
  \AttributeTok{burned\_candidates =}\NormalTok{ refined\_candidates\_path,}
  \AttributeTok{config =}\NormalTok{ cfg,}
  \AttributeTok{keep\_pool =} \ConstantTok{NULL}\NormalTok{,}
  \AttributeTok{write\_outputs =} \ConstantTok{TRUE}\NormalTok{,}
  \AttributeTok{overwrite =} \ConstantTok{TRUE}
\NormalTok{)}

\NormalTok{scored}\SpecialCharTok{$}\NormalTok{internal\_decisions\_path}
\NormalTok{scored}\SpecialCharTok{$}\NormalTok{keep\_pool\_summary}

\FunctionTok{table}\NormalTok{(scored}\SpecialCharTok{$}\NormalTok{internal\_decisions}\SpecialCharTok{$}\NormalTok{class\_final, }\AttributeTok{useNA =} \StringTok{"ifany"}\NormalTok{)}
\FunctionTok{table}\NormalTok{(scored}\SpecialCharTok{$}\NormalTok{internal\_decisions}\SpecialCharTok{$}\NormalTok{filter\_3, }\AttributeTok{useNA =} \StringTok{"ifany"}\NormalTok{)}
\end{Highlighting}
\end{Shaded}

\subsubsection{Validation}\label{validation}

The keep polygons from the scored output can be compared against an
external reference burned-area dataset using
\texttt{validate\_fire\_maps()}.

\begin{Shaded}
\begin{Highlighting}[]
\NormalTok{validation\_dir }\OtherTok{\textless{}{-}} \FunctionTok{file.path}\NormalTok{(cfg}\SpecialCharTok{$}\NormalTok{output\_routes}\SpecialCharTok{$}\NormalTok{base, }\StringTok{"VALIDATION\_KEEP\_ONLY"}\NormalTok{)}
\FunctionTok{dir.create}\NormalTok{(validation\_dir, }\AttributeTok{recursive =} \ConstantTok{TRUE}\NormalTok{, }\AttributeTok{showWarnings =} \ConstantTok{FALSE}\NormalTok{)}

\NormalTok{keep\_path }\OtherTok{\textless{}{-}} \FunctionTok{file.path}\NormalTok{(validation\_dir, }\StringTok{"keep\_only.gpkg"}\NormalTok{)}

\NormalTok{keep\_only }\OtherTok{\textless{}{-}}\NormalTok{ scored}\SpecialCharTok{$}\NormalTok{internal\_decisions[}
\NormalTok{  scored}\SpecialCharTok{$}\NormalTok{internal\_decisions}\SpecialCharTok{$}\NormalTok{class\_final }\SpecialCharTok{==} \StringTok{"keep"}\NormalTok{,}
\NormalTok{]}

\NormalTok{keep\_only }\OtherTok{\textless{}{-}}\NormalTok{ sf}\SpecialCharTok{::}\FunctionTok{st\_zm}\NormalTok{(keep\_only, }\AttributeTok{drop =} \ConstantTok{TRUE}\NormalTok{, }\AttributeTok{what =} \StringTok{"ZM"}\NormalTok{)}
\NormalTok{keep\_only }\OtherTok{\textless{}{-}}\NormalTok{ sf}\SpecialCharTok{::}\FunctionTok{st\_make\_valid}\NormalTok{(keep\_only)}
\NormalTok{keep\_only }\OtherTok{\textless{}{-}}\NormalTok{ keep\_only[}\SpecialCharTok{!}\NormalTok{sf}\SpecialCharTok{::}\FunctionTok{st\_is\_empty}\NormalTok{(keep\_only), , drop }\OtherTok{=} \ConstantTok{FALSE}\NormalTok{]}

\ControlFlowTok{if}\NormalTok{ (}\FunctionTok{nrow}\NormalTok{(keep\_only) }\SpecialCharTok{==} \DecValTok{0}\NormalTok{) \{}
  \FunctionTok{stop}\NormalTok{(}\StringTok{"No polygons with class\_final == \textquotesingle{}keep\textquotesingle{} were found. Validation cannot continue."}\NormalTok{)}
\NormalTok{\}}

\ControlFlowTok{if}\NormalTok{ (}\FunctionTok{file.exists}\NormalTok{(keep\_path)) }\FunctionTok{unlink}\NormalTok{(keep\_path, }\AttributeTok{force =} \ConstantTok{TRUE}\NormalTok{)}

\NormalTok{sf}\SpecialCharTok{::}\FunctionTok{st\_write}\NormalTok{(}
\NormalTok{  keep\_only,}
\NormalTok{  keep\_path,}
  \AttributeTok{layer =} \StringTok{"keep\_only"}\NormalTok{,}
  \AttributeTok{quiet =} \ConstantTok{TRUE}
\NormalTok{)}
\end{Highlighting}
\end{Shaded}

\begin{Shaded}
\begin{Highlighting}[]
\NormalTok{validation\_mask\_path }\OtherTok{\textless{}{-}} \FunctionTok{file.path}\NormalTok{(}
\NormalTok{  validation\_dir,}
  \StringTok{"support"}\NormalTok{,}
  \StringTok{"validation\_mask\_geometry\_only.gpkg"}
\NormalTok{)}

\FunctionTok{dir.create}\NormalTok{(}\FunctionTok{dirname}\NormalTok{(validation\_mask\_path), }\AttributeTok{recursive =} \ConstantTok{TRUE}\NormalTok{, }\AttributeTok{showWarnings =} \ConstantTok{FALSE}\NormalTok{)}

\NormalTok{validation\_mask\_shapefile }\OtherTok{\textless{}{-}} \ControlFlowTok{if}\NormalTok{ (target\_year }\SpecialCharTok{\textless{}} \DecValTok{2000}\NormalTok{) \{}
  \FunctionTok{file.path}\NormalTok{(}
\NormalTok{    base\_dir,}
    \StringTok{"1\_DATA/Mask\_StudyArea"}\NormalTok{,}
    \FunctionTok{paste0}\NormalTok{(}\StringTok{"mask\_3035\_"}\NormalTok{, target\_year, }\StringTok{".shp"}\NormalTok{)}
\NormalTok{  )}
\NormalTok{\} }\ControlFlowTok{else}\NormalTok{ \{}
  \FunctionTok{file.path}\NormalTok{(}
\NormalTok{    base\_dir,}
    \StringTok{"1\_DATA/Mask\_StudyArea"}\NormalTok{,}
    \StringTok{"mask\_Peninsula\_3035.shp"}
\NormalTok{  )}
\NormalTok{\}}

\NormalTok{mask\_sf }\OtherTok{\textless{}{-}}\NormalTok{ sf}\SpecialCharTok{::}\FunctionTok{st\_read}\NormalTok{(validation\_mask\_shapefile, }\AttributeTok{quiet =} \ConstantTok{TRUE}\NormalTok{) }\SpecialCharTok{|\textgreater{}}
\NormalTok{  sf}\SpecialCharTok{::}\FunctionTok{st\_zm}\NormalTok{(}\AttributeTok{drop =} \ConstantTok{TRUE}\NormalTok{, }\AttributeTok{what =} \StringTok{"ZM"}\NormalTok{) }\SpecialCharTok{|\textgreater{}}
\NormalTok{  sf}\SpecialCharTok{::}\FunctionTok{st\_make\_valid}\NormalTok{()}

\NormalTok{mask\_sf }\OtherTok{\textless{}{-}}\NormalTok{ mask\_sf[}\SpecialCharTok{!}\NormalTok{sf}\SpecialCharTok{::}\FunctionTok{st\_is\_empty}\NormalTok{(mask\_sf), }\DecValTok{0}\NormalTok{, drop }\OtherTok{=} \ConstantTok{FALSE}\NormalTok{]}
\NormalTok{mask\_sf}\SpecialCharTok{$}\NormalTok{mask\_id }\OtherTok{\textless{}{-}} \FunctionTok{seq\_len}\NormalTok{(}\FunctionTok{nrow}\NormalTok{(mask\_sf))}

\ControlFlowTok{if}\NormalTok{ (}\FunctionTok{nrow}\NormalTok{(mask\_sf) }\SpecialCharTok{==} \DecValTok{0}\NormalTok{) \{}
  \FunctionTok{stop}\NormalTok{(}\StringTok{"The validation mask is empty after geometry cleaning."}\NormalTok{)}
\NormalTok{\}}

\ControlFlowTok{if}\NormalTok{ (}\FunctionTok{file.exists}\NormalTok{(validation\_mask\_path)) }\FunctionTok{unlink}\NormalTok{(validation\_mask\_path, }\AttributeTok{force =} \ConstantTok{TRUE}\NormalTok{)}

\NormalTok{sf}\SpecialCharTok{::}\FunctionTok{st\_write}\NormalTok{(}
\NormalTok{  mask\_sf,}
\NormalTok{  validation\_mask\_path,}
  \AttributeTok{layer =} \StringTok{"mask"}\NormalTok{,}
  \AttributeTok{quiet =} \ConstantTok{TRUE}
\NormalTok{)}
\end{Highlighting}
\end{Shaded}

\begin{Shaded}
\begin{Highlighting}[]
\NormalTok{strata\_raster }\OtherTok{\textless{}{-}} \FunctionTok{file.path}\NormalTok{(}
\NormalTok{  base\_dir,}
  \StringTok{"1\_DATA/Corine\_Masks/STRATA"}\NormalTok{,}
  \FunctionTok{paste0}\NormalTok{(}\StringTok{"strata\_CLC\_"}\NormalTok{, corine\_year, }\StringTok{"\_res30.tif"}\NormalTok{)}
\NormalTok{)}

\NormalTok{strata\_lut }\OtherTok{\textless{}{-}} \FunctionTok{file.path}\NormalTok{(}
\NormalTok{  base\_dir,}
  \StringTok{"1\_DATA/Corine\_Masks/LUT/lut\_full\_strata8\_v1.csv"}
\NormalTok{)}
\end{Highlighting}
\end{Shaded}

\begin{Shaded}
\begin{Highlighting}[]
\NormalTok{val }\OtherTok{\textless{}{-}} \FunctionTok{validate\_fire\_maps}\NormalTok{(}
  \AttributeTok{input\_shapefile =}\NormalTok{ keep\_path,}
  \AttributeTok{ref\_shapefile =}\NormalTok{ reference\_burned\_map,}
  \AttributeTok{mask\_shapefile =}\NormalTok{ validation\_mask\_path,}
  \AttributeTok{burnable\_raster =}\NormalTok{ burnable\_mask,}
  \AttributeTok{year\_target =}\NormalTok{ target\_year,}
  \AttributeTok{validation\_dir =}\NormalTok{ validation\_dir,}
  \AttributeTok{force\_reprocess\_ref =} \ConstantTok{FALSE}\NormalTok{,}
  \AttributeTok{force\_reprocess\_pred =} \ConstantTok{FALSE}\NormalTok{,}
  \AttributeTok{metrics\_type =} \StringTok{"all"}\NormalTok{,}
  \AttributeTok{dissolve\_ref\_by =} \StringTok{"id"}\NormalTok{,}
  \AttributeTok{dissolve\_input\_by =} \ConstantTok{NULL}\NormalTok{,}
  \AttributeTok{strata\_raster =}\NormalTok{ strata\_raster,}
  \AttributeTok{strata\_lut =}\NormalTok{ strata\_lut,}
  \AttributeTok{observability\_raster =}\NormalTok{ change\_index,}
  \AttributeTok{ref\_end\_doy\_col =} \StringTok{"end\_doy"}\NormalTok{,}
  \AttributeTok{ref\_start\_doy\_col =} \StringTok{"start\_doy"}
\NormalTok{)}
\end{Highlighting}
\end{Shaded}

\begin{tcolorbox}[enhanced jigsaw, arc=.35mm, bottomrule=.15mm, bottomtitle=1mm, breakable, colback=white, colbacktitle=quarto-callout-note-color!10!white, colframe=quarto-callout-note-color-frame, coltitle=black, left=2mm, leftrule=.75mm, opacityback=0, opacitybacktitle=0.6, rightrule=.15mm, title=\textcolor{quarto-callout-note-color}{\faInfo}\hspace{0.5em}{Note}, titlerule=0mm, toprule=.15mm, toptitle=1mm]

This subsection is currently a draft. Later revisions will execute these
steps interactively and discuss the intermediate outputs.

\end{tcolorbox}


\backmatter


\end{document}
